\section{Practical Use and Limitations}
\paragraph{Where the method has a concrete use.}
IC-HACT is a candidate-verification component for a coding or scientific-programming harness. It is useful when workers repeatedly propose changes, the system can execute a meaningful acceptance suite, and users need small, accurate progress reports. Its immediate value is earlier genuine failure evidence, compact verification-state delivery, and explicit rejection of stale or incomplete pass claims. It is not a replacement for a capable coding model, a good task decomposition, or a carefully designed acceptance contract.

A practical integration keeps source-edit ownership in the existing task DAG. Each candidate enters a read-only or privately copied verification stage. The advisory model proposes an order; the trusted wrapper validates that the order covers the approved registry exactly. Results feed both the certificate tree and a durable evidence store. A failure returns a short structured finding to the relevant worker. A complete pass produces an identity-bound local ticket. An outer integration service must then atomically ensure that the checked candidate is the candidate being merged; this final external transaction is not implemented by our Python lock.

\paragraph{Context discipline.}
The certificate tree eliminates the need for a model to reconstruct verification state from full conversations. A root can expose counts, failures, unknown obligations, and evidence identifiers without reading every descendant's logs. However, difficult design questions may still require broad reasoning. The 3072-byte bound applies to a verification packet, not to all reasoning contexts. Tree depth, number of refreshes, raw log size, and operator retrieval remain separate costs. A deep learned tree can worsen progress latency despite respecting the local packet cap.

\paragraph{No causal overclaim.}
A mutation is an intervention, but an intervention dataset is not automatically a complete causal model. Our graph is an empirical, source-conditioned failure relation under frozen fixtures. It can miss indirect dependencies, new control paths, dynamic imports, environmental effects, and interactions among simultaneous edits. The method's safety argument deliberately does not require that graph to be complete. Establishing broader causal invariance would require richer interventions and identifiable assumptions, not a change in terminology.

\paragraph{Threats to empirical validity.}
The study contains two upstream libraries and one related prototype, selected test scopes, synthetic source mutations rather than a representative history of human or LLM repairs, few source-file clusters, and one host. Mutation operators can create unrealistic faults or equivalent changes. Several training files provide very little failure evidence. The confirmation set tests new sites in known projects and mostly known files; it is not an unseen-project transfer study. Interpreting bootstrap intervals outside those boundaries would be unwarranted.

Actual ordering can interact with shared fixtures or order-dependent tests, which is why we compare direct timing-run verdicts with their full-oracle verdicts rather than assuming equivalence. Timeout cases remain unresolved. No experiment evaluates hosted LLM calls, total token consumption, autonomous patch quality, sustained thousands-agent operation, cross-machine scheduling, or high-availability failover. The protocol result is conditional on a trusted checker and accurately identified execution inputs; local source copying alone is not protection against malicious code or external data changes.

\paragraph{What a larger conference evaluation still needs.}
The manuscript supplies a precise research object, proofs, executable controls, negative results, and a reproducible artifact. It does not itself establish acceptance at a leading venue. A stronger external-validity claim requires independent repositories with real revision histories, standardized agent-generated patch tasks, comparison with fully implemented modern test-selection/prioritization systems, prospective workload drift, and end-to-end integrations that account for preparation, training, verification, and deployment costs. These are explicitly outstanding evaluations, not results deferred to an unavailable appendix.

\section{Conclusion}
A certificate layout optimized for what becomes stale can be wrong for how evidence becomes available. IC-HACT makes that distinction explicit: the completion-frontier distribution, including first-failure stopping and initial construction, is the relevant object for progressive verification traffic. A small counterexample proves the information gap; an exact inherited-DP compiler realizes the new objective. Controlled source interventions can help prioritize feedback, but learned structure is kept outside the acceptance predicate. The executable study demonstrates both useful effects and limits, including conditional-layout overfitting and startup-dominated latency. The practical design principle is to let learning optimize the route to evidence, while a complete, current, independently enforced contract determines whether the system may claim success.
